% Section 15: Confluence and decidable derivability
\label{sec:confluence-body}

Normalization in \S\ref{sec:normalization} guarantees that a normal form exists.  It does
not guarantee that the normal form is unique, and uniqueness is what a decision procedure
needs.  This section supplies the critical-pair analysis that yields confluence, then the
normal/neutral characterization that yields the subformula property, and finally the
focused decision procedure that follows.  It closes by describing precisely what finite
decidability does \emph{not} give.

\subsection{Critical pairs}

Because $\mathrm{TND}_0$ is linear and ordered, the only way two redexes can interfere is
that their ledgers overlap.  This makes the critical-pair analysis finite and explicit.

\begin{definition}[Overlapping redexes]\label{def:overlapping-redexes}
Two redexes in a $\mathrm{TND}_0$ proof tree \emph{overlap} when their ledgers share an
antecedent occurrence.  Since ledgers are consumed linearly
(Definition~\ref{def:tnd0}), a shared occurrence means one redex lies on the path from the
other to its discharged assumption.
\end{definition}

\begin{proposition}[Four-way path classification]\label{prop:four-paths}
For an overlapping pair of redexes, exactly one of four configurations occurs:
\begin{enumerate}
\item one redex strictly above the other, on the discharged spine;
\item one redex strictly above the other, off the discharged spine;
\item the two redexes are adjacent along one ledger block;
\item the two redexes are disjoint but order-adjacent in the ledger.
\end{enumerate}
\end{proposition}

\begin{proof}
By Definition~\ref{def:overlapping-redexes} the two redexes lie on one spine, and the
relative position of the shared occurrence in their ledger blocks gives the four cases.
Linearity excludes any configuration in which a ledger block is shared without one redex
containing the other.
\end{proof}

\begin{lemma}[Substitution and graft equations]\label{lem:substitution-graft}
In configurations (1) and (3) the two reduction orders are related by the ledger splice of
Theorem~\ref{thm:substitution-splice}; in configurations (2) and (4) they commute exactly.
\end{lemma}

\begin{proof}
Configurations (2) and (4) act on disjoint ledger blocks, so the root transports of
Theorem~\ref{thm:root-transport} commute by inspection.  Configurations (1) and (3) nest
one substitution inside the other, and the splice law identifies the two composites; the
only thing to check is that the ledger blocks compose in the same order in both, which is
the content of freshness.
\end{proof}

\subsection{The one-step strong diamond}

\begin{theorem}[One-step strong diamond]\label{thm:strong-diamond}
If $P\leadsto P_1$ and $P\leadsto P_2$ with $P_1\neq P_2$, then there is $P_3$ with
$P_1\leadsto P_3$ and $P_2\leadsto P_3$, and the two-step reductions can be chosen to have
equal length.
\end{theorem}

\begin{proof}
By Lemma~\ref{lem:substitution-graft}, each configuration of
Proposition~\ref{prop:four-paths} either commutes or is joined after one further step.  The
equal-length clause follows because both reductions remove the same number of detours in
configurations (2) and (4), and differ by a single commuting step in (1) and (3).
\end{proof}

\begin{corollary}[Local and global confluence]\label{cor:confluence}
The reduction relation of $\mathrm{TND}_0$ is locally confluent and, being strongly
normalizing (Theorem~\ref{thm:strong-normalization}), globally confluent by Newman's
lemma.  Consequently every proof has a unique normal form.
\end{corollary}

\begin{remark}
Corollary~\ref{cor:confluence} is stated for a \emph{fixed} proof tree.  It does not say
that two distinct proofs of the same sequent share a normal form; the fragment is
incomplete (Theorem~\ref{thm:incompleteness}), so distinct derivations of one sequent need
not be related at all.
\end{remark}

\subsection{Confluence is not coherence}

\begin{nogo}[Confluence does not give a 2-cell structure]\label{nogo:confluence-coherence}
The fact that reductions commute is not the fact that they are coherent.  There is a
four-level ladder
\[
  \text{proof endpoint}\ \subsetneq\ \text{composite boundary transport}\ \subsetneq\
  \text{history}\ \subsetneq\ \text{history 2-cell},
\]
and confluence at the first level determines nothing at the fourth.  Reading
Corollary~\ref{cor:confluence} as a coherence theorem is therefore invalid.
\end{nogo}

\begin{proof}
Levels 1 and 2 are quotients of level 3, and level 4 records relations among histories
that are collapsed by the quotient.  Confluence is a statement about the image at level 1;
a nontrivial 2-cell lives in the kernel of the forgetful map from level 3 to level 1, and
that kernel is not determined by level 1.
\end{proof}

\begin{remark}
No-go~\ref{nogo:confluence-coherence} is the proof-theoretic counterpart of
Proposition~\ref{prop:lossy-quotient}.  In both cases the error is the same: taking a
quotient for the object.
\end{remark}

\subsection{Normal and neutral forms}

\begin{definition}[Normal and neutral]\label{def:normal-neutral}
A derivation is \emph{normal} if it contains no principal-implication detour.  It is
\emph{neutral} if it ends in an elimination whose major premise is a hypothesis.  The
standard characterization is: normal derivations end either in an introduction or are
neutral, and neutral derivations are closed under the eliminations of the fragment.
\end{definition}

\begin{theorem}[Subformula property]\label{thm:subformula}
Every formula occurring in a normal $\mathrm{TND}_0$ derivation is a subformula of the end
sequent or of a discharged hypothesis.
\end{theorem}

\begin{proof}
Induction on the normal derivation using Definition~\ref{def:normal-neutral}.  An
introduction rule's premises contain only subformulas of its conclusion; an elimination
whose major premise is neutral has a minor premise that is a subformula of the major
premise's type, and the absence of detours means no formula is introduced only to be
eliminated.
\end{proof}

\begin{proposition}[Leftmost-head canonicity]\label{prop:leftmost-head}
In a normal derivation of an implicational sequent, the construction of the proof is
determined by the ordered context together with the leftmost hypothesis of the end
sequent's antecedent block.  Hence the leftmost-head strategy is complete for the
fragment.
\end{proposition}

\begin{proof}
By Theorem~\ref{thm:subformula} every formula of a normal derivation is a subformula of the
goal, and by linearity no antecedent is reused, so at each stage the possible rules are
determined by the leftmost unconsumed antecedent and the goal.  The choices are therefore
forced up to the order in which independent threads are processed.
\end{proof}

\subsection{Decidable derivability}

\begin{theorem}[Focused decision procedure]\label{thm:decidability}
Derivability in $\mathrm{TND}_0$ is decidable.  There is a terminating, sound, and complete
procedure: search normal derivations in leftmost-head order, using
Theorem~\ref{thm:subformula} to bound the formula set, Definition~\ref{def:normal-neutral}
to alternate focused and unfocused phases, and Proposition~\ref{prop:leftmost-head} to fix
the branching.
\end{theorem}

\begin{proof}
Termination: by Theorem~\ref{thm:subformula} only finitely many formulas may occur, and
linearity bounds the number of introductions by the number of antecedents.  Soundness: each
step of the procedure builds a derivation, so a found proof is a proof.  Completeness: by
Corollary~\ref{cor:confluence} a normal form exists if any derivation does, and by
Proposition~\ref{prop:leftmost-head} the procedure explores a branch containing it.
\end{proof}

\begin{proposition}[Positioning of the skeleton]\label{prop:skeleton-position}
The $\mathrm{TND}_0$ skeleton is the associative, empty-antecedent, right-residual fragment
of ordered linear logic.  Under this reading, the erasure of audited nominal data to a
skeleton is a functor, and lifting back is by fresh, noncanonical decoration; long forms
are not identified with their eta-expansions, and multiple focused proofs of one sequent
are admitted.
\end{proposition}

\begin{proof}[Status]
The fragment identification is a syntactic matching of rule shapes.
Proposition~\ref{prop:leftmost-head} and Theorem~\ref{thm:decidability} are the substantive
content; the erasure/lift discipline is a definitional choice, recorded so that the
decision procedure's non-canonicity of decoration is visible rather than implicit.
\end{proof}

\subsection{What finite decidability does not give}

\begin{warning}[No compactification from finite search]\label{warn:no-compactification}
Theorem~\ref{thm:decidability} is a statement about a finite fragment.  It does not
produce, and must not be read as producing, a halting measure, a limit object, or any
$\Omega$-type boundary.  Neither finite confluence
(Corollary~\ref{cor:confluence}) nor finite focused search
(Theorem~\ref{thm:decidability}) yields such an object; they are finite statements and the
boundary is not their limit.
\end{warning}

\begin{remark}
Warning~\ref{warn:no-compactification} is the sharpest instance of the discipline that
organises this paper.  A decision procedure terminates; that is not the same as exhibiting
a limit of a non-terminating process, and conflating the two would convert a finite theorem
into an infinite claim.
\end{remark}

\subsection{Summary}

\begin{enumerate}
\item The critical-pair analysis is finite, with four configurations
  (Proposition~\ref{prop:four-paths}, Lemma~\ref{lem:substitution-graft}).
\item Reduction satisfies a one-step strong diamond and is confluent, so normal forms are
  unique per proof (Theorem~\ref{thm:strong-diamond}, Corollary~\ref{cor:confluence}).
\item Confluence is not coherence (No-go~\ref{nogo:confluence-coherence}).
\item The subformula property plus leftmost-head canonicity give a terminating, sound, and
  complete decision procedure (Theorems~\ref{thm:subformula}, \ref{thm:decidability},
  Proposition~\ref{prop:leftmost-head}).
\item Decidability yields no $\Omega$-type boundary
  (Warning~\ref{warn:no-compactification}).
\end{enumerate}
